\documentclass{article}
\usepackage[utf8]{inputenc}
\usepackage[absolute]{textpos}
\usepackage[default]{raleway}
\usepackage{titlesec, comment, tabularx, makecell, listings, array, setspace, geometry, graphicx, xcolor, xparse, fancyvrb, relsize, fancyhdr, booktabs, hyperref}
\usepackage{colortbl}
\usepackage{float}
%\geometry{a4paper, left=2cm, right=2cm, top=2cm, bottom=2.5cm}
\renewcommand{\headrulewidth}{0pt}

% Definisci uno stile per i comandi git
\definecolor{light-gray}{gray}{0.92}

\lstdefinestyle{code}{
    frame=single,
    framesep=1mm,
    rulecolor=\color{light-gray},
    backgroundcolor=\color{light-gray},
    basicstyle=\ttfamily,
} 

% ----------------------------- Definizione tabella ---------------------------

\newcolumntype{C}[1]{>{\centering\arraybackslash}m{#1}}

%\setcellgapes{2ex} % Imposta l'altezza dell'header (2ex)


% ------------------------------Metadati indice --------------------------------
\title{\textbf{\fontsize{28}{6}\selectfont Indice}}
\author{\fontsize{14}{6}\selectfont ByteOps}
\date{Dicembre 8, 2023}

% -----------------------------Creazione footer --------------------------------

\pagestyle{fancy}
\fancyhf{}
\renewcommand{\footrulewidth}{0.4pt}
\lfoot{
    \parbox[c]{2cm}{\includegraphics[width=2cm]{../../../Images/logo.png}}
    \textcolor[RGB]{120, 120, 120}{$\cdot$ Verbale Interno}
}
\rfoot{\thepage}

% --------------------------Modifica formato hyperlinks ------------------------

\hypersetup{
    colorlinks=true,
    linkcolor=black,
    filecolor=black,      
    pdftitle={Verbale Interno 04/12/2023},  %inserisci data verbale
    pdfpagemode=FullScreen,
}

% ------------------------------- Valore sotto-paragrafi indice --------------------------------------

\setcounter{secnumdepth}{4}
\setcounter{tocdepth}{4}

\titleformat{\section}
{\normalfont\huge\bfseries}{\thesection}{0.2cm}{}
\titlespacing*{\paragraph}{0pt}{0.5cm}{0.1cm}

\titleformat{\subsection}
{\normalfont\Large\bfseries}{\thesubsection}{0.2cm}{}
\titlespacing*{\paragraph}{0pt}{0.5cm}{0.1cm}

\titleformat{\subsubsection}
{\normalfont\large\bfseries}{\thesubsubsection}{0.2cm}{}
\titlespacing*{\paragraph}{0pt}{0.5cm}{0.1cm}

\titleformat{\paragraph}
{\normalfont\normalsize\bfseries}{\theparagraph}{0.2cm}{}
\titlespacing*{\paragraph}{0pt}{0.5cm}{0.1cm}

% ------------------------------- Front Page ---------------------------------------

\begin{document}

% --------------------------Aggiunta firma finale ------------------------
\begin{textblock*}{\textwidth}(0.85\textwidth, 1.16\textheight)
    Il responsabile: Riccardo Smanio
\end{textblock*}
% ------------------------------------------------------------------------

\pagestyle{fancy}
\begin{center}
\includegraphics[width = 0.7\textwidth]{../../../Images/logo.png} \\
\vspace{0.2cm}
\textcolor[RGB]{60, 60, 60}{\textit{ByteOps.swe@gmail.com}} \\
\vspace{1cm}
\fontsize{16}{6}\selectfont Verbale Interno $\cdot$ Data: 08/12/2023 \\
\vspace{0.5cm}
\end{center}

\section*{Informazioni documento}
\def\arraystretch{1.2}
\begin{tabular}{>{\raggedleft\arraybackslash}p{0.2\textwidth}|>{\raggedright\arraybackslash}p{0.6\textwidth}c}
\hline
\addlinespace
\textbf{Luogo} & Discord \vspace{10pt} \\
\textbf{Orario} & 10:30 - 11:30 \vspace{10pt} \\
\textbf{Redattore} & A. Barutta \vspace{10pt} \\
\textbf{Verificatore} & N. Preto \vspace{10pt} \\
\textbf{Amministratore} & D. Diotto \vspace{10pt} \\
\textbf{Destinatari} & T. Vardanega \\ & R. Cardin \vspace{10pt} \\
\textbf{Partecipanti} & A. Barutta \\ & E. Hysa \\ & R. Smanio \\ & D. Diotto \\ & F. Pozza \\ & L. Skenderi \\ & N. Preto \vspace{10pt} \\
\end{tabular}
\pagebreak 

% ------------------------- Changelog ----------------------------

\section*{Registro delle modifiche}

\begin{tabular}{|C{2.5cm}|C{2.5cm}|C{2.5cm}|C{2.5cm}|C{2.5cm}|}
    \hline
    \textbf{Versione} & \textbf{Data} & \textbf{Autore} & \textbf{Verificatore} & \textbf{Dettaglio} \\
    \hline \hline
    0.0.1 & 09/12/2023 & Andrea Barutta & N. Preto & Redazione completa. \\
    \hline
\end{tabular}
\pagebreak

% ------------------------- Generazione automatica indice ----------------------
\setstretch{1.5}
\maketitle
\thispagestyle{fancy}
\tableofcontents
\setstretch{1.2}
\pagebreak

% ------------------------ INIZIO DOCUMENTO ----------------------
\flushleft

\section{Revisione del periodo precedente}
Nella fase iniziale della riunione, è stata condotta un'analisi approfondita sullo stato delle attività in corso, rilevando che procedono conformemente alle previsioni, senza manifestare alcuna problematica evidente. Di conseguenza, non sono emerse proposte di modifica ai processi né sono state sollevate discussioni rilevanti.  
\section{Ordine del giorno}
    \subsection{Cruscotto di valutazione della qualità - Piano di qualifica}
    Dopo aver inizialmente redatto le metriche di qualità di processo e prodotto, è essenziale condurre una revisione approfondita e sviluppare un cruscotto di valutazione per tracciare il valore di queste metriche durante i diversi periodi del progetto. Questo cruscotto consentirà di visualizzare l'andamento e facilitare valutazioni migliorative. La decisione presa è stata di creare un documento condiviso su Google Fogli, dove verranno registrati i valori delle metriche per ogni periodo e saranno generati grafici illustrativi degli andamenti.  
    \subsection{Revisione completa - Norme di progetto}
    Dopo aver completato la stesura del documento "Norme di progetto", si è reputato indispensabile condurre una revisione completa al fine di verificare la coerenza e la consistenza delle diverse sezioni redatte, garantendo così che il file sia pronto per la revisione RTB.  
    \subsection{Ottimizzazione POC}
    In seguito al feedback positivo ricevuto durante il SAL del 06/12/2023, si è concordato con il proponente di considerare complete le funzionalità richieste per il POC. Di conseguenza, il focus dello sviluppo sarà ora sulla risoluzione di piccoli bug e sulla creazione di una versione stabile da presentare durante la revisione RTB.  
    \subsection{Feedback nuova versione dei Casi d'Uso}
    In seguito all'incontro con il Prof. Cardin, si è proceduto alla ridefinizione dei casi d'uso e, nel seguente meeting, sono stati sottoposti all'analisi dell'intero team al fine di raccogliere un feedback da parte di ogni componente.\\
    Considerato il riscontro positivo, l'obiettivo per il prossimo sprint è la creazione di una nuova versione del documento "Analisi dei Requisiti" da sottoporre nuovamente alla valutazione del professore.
    \subsection{Rotazione ruoli} 

    \begin{table}[H]
        \centering
        \begin{tabular}{|c|c|} 
            \hline
            \textbf{Ruolo} & \textbf{Nome Cognome} \\
            \hline \hline
            Responsabile (Re) & R. Smanio\\ 
            \hline
            Amministratore (Am) & D. Diotto\\ 
            \hline
            Analista (An) & L. Skenderi \\
            \hline
            Verificatore (Ve) & \makecell{N. Preto\\A. Barutta } \\
            \hline
            Programmatore (Pr) & \makecell{F. Pozza\\E.Hysa } \\
            \hline
        \end{tabular}
        \caption{Tabella che rappresenta i membri per ogni ruolo}
    \end{table}

\section{Attività da svolgere}
    \begin{table}[H]
        \centering
        \begin{tabular}{|C{7cm}|C{1,5cm}|C{3cm}|}
            \hline
            \textbf{Titolo} & \textbf{\# Issue} & \textbf{Verificatore}\\
            \hline\hline           
            Piano di qualifica - cruscotto valutazione e grafici per il miglioramento  & 55 & A. Barutta\\
            \hline
            Revisione finale Norme di Progetto & 56 & A. Barutta\\
            \hline
            Ottimizzazione POC: risoluzione bug e revisione RTB & 57 & N. Preto\\
            \hline
            Redazione verbale interno 08/12/2023 & 58 & N. Preto\\
            \hline
            Analisi dei requisiti - redazione nuova versione dei casi d'uso & 59 & A. Barutta\\
            \hline
        \end{tabular}
        \caption{Tabella conentente le attività da svolgere nel prossimo sprint}
    \end{table}

\end{document}